\uuid{H9U4}
\titre{ Règle de mise à jour en descente de gradient }
\niveau{L3}
\module{Analyse de données}
\chapitre{Réseaux de neurones}
\sousChapitre{Descente de gradient, mise à jour des paramètres}
\theme{Signe du gradient, minimisation du coût}
\auteur{Maxime Nguyen}
\datecreate{2026-03-19}
\organisation{AMSCC}
\difficulte{1}

\contenu{
\texte{
  On cherche à minimiser un coût $J$ par descente de gradient.
  On compare trois règles de mise à jour :
  \begin{itemize}
    \item[A.] \texttt{w = w + lr * gradient}
    \item[B.] \texttt{w = w - lr * gradient}
    \item[C.] \texttt{w = w * lr * gradient}
  \end{itemize}
}
\begin{enumerate}

\item
  \question{Laquelle est correcte pour minimiser $J$ ? Justifier les erreurs des deux autres.}
  \reponse{
    \textbf{B} est correcte : la descente de gradient soustrait le gradient pour aller dans la direction de plus grande décroissance.
    \textbf{A} effectue une \emph{montée} de gradient, ce qui maximise $J$ au lieu de le minimiser.
    \textbf{C} est une opération sans sens mathématique dans ce contexte : elle modifie l'ordre de grandeur de $w$
    de façon multiplicative et ne correspond à aucun algorithme d'optimisation standard.
  }

\end{enumerate}
}
